\section*{Leitidee}

Funktionen beschreiben, wie Gr��en voneinander abh�ngen. Neben ihrem
globalen Verlauf ist dabei h�ufig entscheidend, was in unmittelbarer
Umgebung eines bestimmten Punktes geschieht: N�hern sich Funktionswerte
einem bestimmten Wert? Ver�ndert sich die Funktion dort stetig? Wie stark
steigt oder f�llt sie?

Die mathematische Grundlage f�r solche lokalen Aussagen bilden
\textbf{Grenzwerte}. Sie erm�glichen es, Ann�herungsprozesse pr�zise zu
beschreiben, und f�hren zu den Begriffen der \textbf{Stetigkeit} und
\textbf{Differenzierbarkeit}.

Der zentrale Gedanke der Differentialrechnung ist die
\textbf{lokale Linearit�t}: In hinreichend kleiner Umgebung l�sst sich eine
differenzierbare Funktion durch ihre Tangente, also durch eine affine
Funktion, approximieren. Die Ableitung beschreibt dabei die lokale
�nderungsrate und damit die Steigung dieser Tangente.

Diese lokale Approximation verbindet die geometrische Vorstellung einer
Tangente mit der Analyse von Ver�nderungen. Sie bildet zugleich die
Grundlage f�r weiterf�hrende Konzepte wie Taylor-Approximationen,
numerische Verfahren und die Beschreibung dynamischer Prozesse durch
Differentialgleichungen.

